\documentclass[../cours_tle.tex]{subfiles}
\begin{document}
\setcounter{chapter}{26}
\chapter{Polynômes, racines et factorisation dans le plan complexe}
\label{ch:27}
\begin{objectif}
Mathématiques expertes. Démontrer ce que le degré permet d'affirmer. Prérequis : chapitres 4 et 25.
\end{objectif}
\section{Une racine produit un facteur}
Un polynôme non nul $P(z)=a_0+a_1z+\cdots+a_nz^n$, $a_n\ne0$, a degré $n$. Une racine est un $a$ tel que $P(a)=0$.
\begin{propriete}
Pour $n\ge1$,
$z^n-a^n=(z-a)\sum_{k=0}^{n-1}z^{n-1-k}a^k$.
Si $P(a)=0$, il existe un polynôme $Q$ tel que $P(z)=(z-a)Q(z)$.
\end{propriete}
\begin{demonstration}
Développer le produit de la première formule : tous les termes intermédiaires s'annulent.
Puis écrire $P(z)-P(a)=\sum_{j=1}^n a_j(z^j-a^j)$ ; chacun des termes est divisible par $z-a$.
\end{demonstration}
\begin{theoreme}
Un polynôme non nul de degré $n$ a au plus $n$ racines distinctes.
\end{theoreme}
\begin{demonstration}
Un polynôme constant non nul n'a pas de racine. Si $P$ de degré $n$ possède une racine $a$, factoriser $P=(z-a)Q$, avec $Q$ de degré $n-1$. Toute racine différente de $a$ est racine de $Q$, puisque $\mathbb C$ est un corps. Conclure par récurrence.
\end{demonstration}
Ainsi un polynôme nul en tous les points a tous ses coefficients nuls. Une infinité de vérifications exactes peut forcer cette conclusion ; quelques valeurs numériques ne suffisent pas sans borne sur le degré.
\section{Second degré à coefficients réels}
Pour $az^2+bz+c=0$, $a\ne0$, poser $\Delta=b^2-4ac$.
La forme canonique donne $(2az+b)^2=\Delta$.
Si $\Delta\ge0$, les solutions réelles usuelles restent valables. Si $\Delta<0$, les racines sont
$(-b\pm i\sqrt{-\Delta})/(2a)$, conjuguées l'une de l'autre.
Si $\Delta=0$, il y a une racine double.
\begin{exemple}
$z^2-2z+5=0$ a pour racines $1\pm2i$.
Pour $P(z)=z^3-z^2+z-1$, $P(1)=0$ et
$P(z)=(z-1)(z^2+1)$, donc les racines sont $1,i,-i$.
\end{exemple}
\section{Racines carrées complexes et coefficients complexes}
Pour résoudre $w^2=a+ib$, poser $w=x+iy$. On obtient
$x^2-y^2=a$, $2xy=b$, et $x^2+y^2=\sqrt{a^2+b^2}$.
Donc $x^2=(|z|+a)/2$, $y^2=(|z|-a)/2$, avec signes choisis pour respecter $2xy=b$.
Pour $3+4i$, les racines carrées sont $\pm(2+i)$.
La forme canonique du second degré reste valable pour des coefficients complexes ; choisir une racine carrée $\delta$ de $\Delta$ donne $(-b\pm\delta)/(2a)$.
\section{Somme, produit et multiplicité}
Si $P=a\prod_{j=1}^n(z-r_j)$, le coefficient de $z^{n-1}$ vaut $-a\sum r_j$ et le terme constant vaut $a(-1)^n\prod r_j$ : choisir les termes appropriés dans le développement. Ce sont les premières formules de Viète.
Une racine de multiplicité $m$ signifie que $(z-a)^m$ divise $P$ mais pas $(z-a)^{m+1}$.
Pour un polynôme à coefficients réels, $P(\overline z)=\overline{P(z)}$ ; les racines non réelles vont donc par paires conjuguées.
Le théorème affirmant qu'un polynôme complexe de degré $n\ge1$ possède exactement $n$ racines avec multiplicité est le théorème fondamental de l'algèbre, admis ici et réservé à une preuve universitaire. La borne « au plus $n$ » ci-dessus ne suffit pas à le démontrer.
\begin{finchapitre}
\section{Utilité et suite du parcours}
Les polynômes relient équations, interpolation, valeurs propres et structures algébriques. La multiplicité prépare l'étude fine des racines et des matrices.
\begin{retenir}
Distinguer racines distinctes et racines comptées avec multiplicité. Toujours exclure le polynôme nul dans une borne fondée sur le degré.
\end{retenir}
\end{finchapitre}
\end{document}
